% === APPENDIX ================================================================
\section{Assumptions and Reproducibility Notes}
\label{subsec:assumptions}

The following are assumptions or design choices rather than measured facts,
listed so a reader can judge how far each result depends on them.

\emph{A1. Yaw inertia.} $I_{z}=\SI{51.1}{\kilo\gram\meter\squared}$ is carried
in the model file and is \emph{not} measured. It enters \eqref{eq:omdot}
directly, so it scales the yaw-rate residual the network is trained on. It can
be discharged by applying a known yaw moment and fitting $I_{z}$ from the
measured $\dot\omega$, or by reading the actor's inertia tensor.

\emph{A2. Centre-of-gravity split.} $l_{f}=\SI{0.738142}{\meter}$,
$l_{r}=\SI{0.795362}{\meter}$ are taken from the controller configuration. The
measured static load split (51.2/48.8) implies a marginally different split than
the geometry does (51.9/48.1), a \SI{1.4}{\percent} disagreement, below the
level at which any conclusion here would change.

\emph{A3--A4. CG height and steering sign.} $h_{\mathrm{cg}}=\SI{0.3}{\meter}$
is carried but not measured. The single-track identification model neglects
load transfer, following the baseline, and its normal loads are static; the
only consumer of $h_{\mathrm{cg}}$ is the longitudinal transfer term of the
friction warm start (Section~\ref{subsec:friction}), where it moves the axle
loads by up to \SI{14}{\percent} at the \SI{0.35}{\g} the fit admits. \texttt{feedback/\allowbreak
steering\_angle} is published in degrees and treated as the front road-wheel
angle, positive left, verified against the front-left wheel joint state to
within \SI{0.5}{\percent}; left and right road wheels are not distinguished.

\emph{A5. Steady-state force recovery.} Equation~\eqref{eq:ssforces} is not a
force measurement; being fully determined by $v_{x}\omega$ it misattributes
transients, and it is applied only to the \emph{synthetic}, quasi-steady
rollout. Its adequacy was checked on measured data against the simulator's own
axle force (correlation \num{0.993}, slope \num{0.94}).

\emph{A6. The residual network is not the plant.} The corrected model is queried
by the rollout at inputs the network has only partially seen. The extrapolation
bound of Section~\ref{subsec:iterative} limits but does not eliminate this: the
sweep still reaches $1.5\times$ the observed 99th percentile of $|\delta|$, a
choice that is empirical rather than derived.

\emph{A7. Simulator determinism and timing.} CARLA runs in synchronous mode with
a \SI{0.033}{\second} physics step; sim time advances faster than wall time, at
a ratio set by the host load and therefore not identical across runs
(between \num{1.28} and \num{1.66} on the four runs of
Section~\ref{sec:results}). All
identification and force-validation data is stamped from the simulation clock,
while the controller node stamps its own metrics on the wall clock; the two are
never mixed inside one metric (Section~\ref{subsec:setup}).

\emph{A8. Software configuration of record.} ROS~2 Humble, Python 3.10, PyTorch
with CUDA on an RTX~2080-class GPU; the bridge is a separate overlay; CARLA
server 0.9.16 / Unreal Engine 4.26; acados 0.5.5 with partial-condensing HPIPM.
The identifier's acquisition timer runs at \SI{50}{\hertz} (the controller runs
at \SI{20}{\hertz}, a \SI{50}{\milli\second} period) with a \SI{30}{\second}
buffer, six co-identification
iterations, 200 training epochs, full-batch Adam at $5\times10^{-4}$, and
$T_{s}=\SI{0.035}{\second}$. \texttt{nn\_train} seeds PyTorch and NumPy at
\num{42}; \texttt{pacejka\_solver.seed} is \texttt{null} in both shipped
parameter sets, so the reported identifications are reproducible in distribution
but not bit for bit, while the sweep of Table~\ref{tab:sweep} is seeded
explicitly and is bit-reproducible.

\emph{A9. Reference trajectory.} Figure~\ref{fig:raceline} shows the raceline
all four runs follow and the lateral demand it places on the plant.

\begin{figure*}[t]
\centering
\includegraphics[width=\textwidth]{raceline.pdf}
\caption{The optimised raceline used for all four
runs~\cite{tumracetraj,heilmeier2020}. (a) Track layout, coloured by lateral
demand, with the marked corner ($R=\SI{18.7}{\meter}$) the tightest on the lap.
(b) Reference speed profile. (c) Lateral demand $v_{x}^{2}|\kappa|$ along the
lap against the flat \SI{5.0}{\meter\per\second\squared} envelope handed to the
optimiser (dashed) and the grip the vehicle has on the static surface (dotted,
$\approx\SI{0.70}{\g}$, measured). No segment exceeds the plant's grip there;
on the decaying surface the dotted line falls through the solid one.}
\label{fig:raceline}
\end{figure*}


\section{Peak-Friction Warm Start on a Synthetic Plant}
\label{subsec:mu}

The friction warm start is evaluated on a synthetic single-track plant with
brush tires, because only a synthetic plant supplies a $\mu$ known exactly and
an excitation level that can be swept independently of it: a swept-sine steering
input for \SI{40}{\second} at \SI{50}{\hertz} with true $\mu=0.70$, the same
effective peak as the static CARLA surface, the steering amplitude setting how
far up the tire curve the manoeuvre reaches (Table~\ref{tab:mu}).

\begin{table}[t]
\caption{Peak friction recovered from a synthetic brush-tire plant (true $\mu=0.70$), against excitation.}
\label{tab:mu}
\centering
\begin{threeparttable}
\small
\begin{tabular}{@{}cccc@{}}
\toprule
$\delta$ amplitude & Grip reached & Bound $\hat\mu$ & Brush fit
\eqref{eq:brush}\\
\midrule
\SI{0.033}{\radian} & \SI{45}{\percent}  & 0.312 & gated off\\
\SI{0.047}{\radian} & \SI{63}{\percent}  & 0.442 & \textbf{0.699}\\
\SI{0.060}{\radian} & \SI{80}{\percent}  & 0.557 & \textbf{0.699}\\
\SI{0.090}{\radian} & \SI{100}{\percent} & 0.700 & \textbf{0.697}\\
\bottomrule
\end{tabular}
\begin{tablenotes}[flushleft]\scriptsize
\item Column~2 is the fraction of the available grip the manoeuvre reached,
$\max|F_{y}|/(\mu F_{z})$; column~3 is the utilisation bound
$\max|F_{y}|/F_{z}$, which is what that fraction leaves as a lower bound on
$\mu$. The brush fit of \eqref{eq:brush}
instead recovers the peak from the curvature of $F_{y}(\alpha)$ and is
therefore flat in excitation.
\item The first row is refused by the absolute gate
$\max|F_{y}|/F_{z}\ge\num{0.40}$ of Section~\ref{subsec:friction}: at
$\mu=0.70$, \SI{45}{\percent} utilisation is only \num{0.31} of $F_{z}$. The
utilisation bound is reported instead.
\end{tablenotes}
\end{threeparttable}
\end{table}

At \SI{63}{\percent} utilisation, a level an ordinary lap reaches, the
utilisation bound is \SI{37}{\percent} low while the brush fit is within
\SI{0.2}{\percent}, and it stays there across true $\mu\in\{0.7,1.05,1.4\}$ on
noise-free data; with odometry and IMU noise injected the error is
grip-dependent, \SI{20}{\percent} at \SI{63}{\percent} utilisation and
\SI{0.4}{\percent} at the peak for $\mu=0.70$. Below the absolute \num{0.40}
floor on $|F_{y}|/F_{z}$ the gate suppresses the fit and reports the bound, so
the failure is visible rather than silent. The gate behaves the same way on the
CARLA runs, rejecting the first buffer of each and returning $\mu$ within
\SI{11}{\percent} of the static plant at an absolute axle utilisation of
\numrange{0.54}{0.68}. That \SI{11}{\percent} bias is the cost of a real
rollout, a reconstructed $I_{z}\dot\omega$ term and a curve whose sustained
demand is only about \num{0.70} of the axle's peak force, and it dominates the
identified peak factor's error on both surfaces. The decaying surface adds a
lag rather than a new error term.

\section{Attribution by Offline Replay}
\label{sec:ablations}

This appendix attributes the closed-loop ordering of
Section~\ref{subsec:summary} to individual settings, at the cost of leaving the
loop. Every number below is an \emph{offline replay} of the identification
stage on the five \SI{30}{\second} buffers each arm collected during its own
run, scored with the quantity the closed-loop tire-force benchmark itself
reports, so these numbers and Table~\ref{tab:comparison}'s are on the same
scale. Section~\ref{subsec:ablations} summarises what follows.

\subsection*{Damping the co-identification loop}
Table~\ref{tab:beta} sweeps $\beta$ of \eqref{eq:relax} at six iterations
throughout; two columns are the argument. At $\beta=\num{1.00}$, the inherited
undamped update, axle force RMSE is \num{70.4} and \SI{62.3}{\newton}, reached
by degrading monotonically from \num{38.9} and \SI{26.3}{\newton} at a single
iteration: the loop amplifies its own extrapolation instead of refining the
estimate. At $\beta=\num{0.00}$ the coefficients never move from the pinned
prior, which is no identification at all, and that is the worst column,
\num{92.5} and \SI{71.8}{\newton}. The minimum sits between them at
$\beta=\num{0.20}$, \num{28.0} and \SI{19.4}{\newton}, a mean of
\SI{23.7}{\newton}, and the identified $BCD$ rises with $\beta$ over the whole
damped range (\num{7.82} at $\beta=0$ to \num{12.75} at $\beta=\num{0.50}$; the
undamped \num{12.42} is no longer ordered by it), so the damping trades how
much of each fit is accepted against how much extrapolation error comes with
it. Six iterations are therefore defensible on this evidence, provided the
update is damped.

\begin{table}[t]
\caption{Relaxation sweep. Five recorded buffers, six iterations, offline replay.}
\label{tab:beta}
\centering
\begin{threeparttable}
\scriptsize
\setlength{\tabcolsep}{3.2pt}
\begin{tabular}{@{}lrrrrrrrrr@{}}
\toprule
$\beta$ & 0.00 & 0.05 & 0.10 & 0.15 & \textbf{0.20} & 0.25 & 0.30 & 0.50 & 1.00\\
\midrule
Front $F_{y}$ RMSE [\si{\newton}] & 92.5 & 68.3 & 48.6 & 33.9 & \textbf{28.0} & 31.2 & 36.4 & 55.6 & 70.4\\
Rear $F_{y}$ RMSE [\si{\newton}] & 71.8 & 49.0 & 31.3 & 19.8 & \textbf{19.4} & 26.0 & 30.2 & 38.4 & 62.3\\
Front $BCD$ [\si{\per\radian}] & 7.82 & 8.82 & 9.66 & 10.36 & 10.88 & 11.32 & 11.65 & 12.75 & 12.42\\
\bottomrule
\end{tabular}
\begin{tablenotes}[flushleft]\scriptsize
\item $\beta=\num{1.00}$ is the inherited undamped update and $\beta=\num{0.00}$ the pinned prior, i.e.\ no identification. \num{0.20} ships.
\end{tablenotes}
\end{threeparttable}
\end{table}

\subsection*{The other settings}
Pinning the peak factor accounts for most of what remains, and the decaying
surface is where its value is easiest to state: an estimator that cannot move
$D$ has nothing to publish when the surface changes. Replayed on five buffers
of a $\mu=\num{1.05}$ run, letting the fit move $D$ gives a mean $|D-\mu|$ of
\num{0.585} and a front-stiffness spread of
\SI{32.9}{\kilo\newton\per\radian}; pinning only the axles whose brush fit was
released reduces those to \num{0.191} and \num{32.3}, and also holding the axles
with no released measurement, the shipped rule, to \num{0.084} and
\SI{18.5}{\kilo\newton\per\radian}. What removes most of the error is refusing
to let an unmeasured $D$ be fitted at all. Normalising the physics-informed
objective is necessary but not sufficient alone: it moves the loss shares from
\SI{99.9996}{\percent} on the dynamics residual to \SI{1.47}{\percent}, yet
applied without the damping it makes the force metric slightly \emph{worse}
(\num{77.1} and \SI{76.2}{\newton} against \num{70.4} and \SI{62.3}{\newton}),
because an objective that now fits the data supplies a sharper trajectory for
the undamped loop to amplify. Two settings deserve a weaker statement than the
closed-loop comparison alone would support: the S4D residual is worth
\SI{21}{\percent} on the force metric at matched settings and a single iteration
(\SI{35.3}{\newton} against the baseline MLP's \SI{44.7}{\newton}), an offline
replay and not a closed-loop result; and the residual \emph{objective} is a
wash, \texttt{physics\_informed} and plain \texttt{mse} being indistinguishable
at \SI{23.7}{\newton} each once the normalisation and the damping are both in
place. We keep it because it costs nothing and is the paper's stated
formulation, not because it buys accuracy. The fast-$\mu$ derate is likewise
not separated: it is active only in \oursarm{}, and on the decaying surface it
sat at the upper end of its $[\num{0.4},\num{1.2}]$ clamp over the last cycles,
so the reference re-clamp of Section~\ref{subsec:resultsclosedloop} is carried
by the published $D$ and not by the derate.

\section{Simulator--ROS~2 Interface}
\label{sec:iface}

\subsection*{Feedback path}
Figure~\ref{fig:bridge} shows the bridge and Table~\ref{tab:topics} lists the
interface consumed by the identifier and its validation.

\begin{figure*}[t]
\centering
\resizebox{\textwidth}{!}{% =============================================================================
% Fig. 3 -- Carla_ASU_Bridge architecture.
%
% Vector TikZ, drawn at the paper's own font sizes. Included by
% sections/method.tex inside a figure* (full text width).
%
% Reading order, left to right: CARLA server -> bridge lifecycle node ->
% ROS 2 topic layer -> autonomy stack.
%   * orange  = the ground-truth path (feedback/tire_forces). It leaves the
%               simulator and reaches only the scoring/validation nodes; it is
%               deliberately NOT wired into the identification loop.
%   * blue    = mpc/update_params, the gated model interface. Drawn inside the
%               autonomy stack because it is a service between two of this
%               workspace's nodes, not part of the bridge's topic layer.
%
% Nodes are stacked with relative positioning (below=... of ...) so that a
% change of wording cannot make two boxes overlap.
% =============================================================================

\begin{tikzpicture}[
  font=\scriptsize,
  >={Stealth[length=1.7mm,width=1.5mm]},
  % --- block styles ---
  colhead/.style  = {font=\scriptsize\bfseries, align=center},
  blk/.style      = {draw, rounded corners=1pt, align=left, inner sep=2.4pt,
                     text width=#1, fill=white, line width=0.35pt},
  srv/.style      = {blk=26mm, fill=black!5},
  brg/.style      = {blk=33mm, fill=black!3},
  top/.style      = {blk=36mm, font=\scriptsize\ttfamily, fill=white},
  gt/.style       = {top, fill=orange!20, draw=orange!80!black, line width=0.7pt},
  cons/.style     = {blk=31mm, fill=black!3},
  note/.style     = {font=\tiny, align=left, inner sep=1.5pt, text width=36mm},
  % --- edge styles ---
  flow/.style     = {->, line width=0.5pt, black!65},
  gtflow/.style   = {->, line width=0.9pt, orange!80!black},
  gateflow/.style = {->, line width=0.9pt, blue!60!black},
  lbl/.style      = {font=\tiny, inner sep=1.2pt, fill=white, midway}
]

\def\gap{3.2mm}   % vertical gap between stacked boxes

% ============================ 1. CARLA server ============================
\node[srv, anchor=north west] (ego) at (0,0)
  {\textbf{Ego vehicle}\\[1pt]
   {\ttfamily vehicle.vehicle.}\\{\ttfamily asurt\_fsai}\\[1pt]
   PhysX wheeled-vehicle solver,\\per-wheel tire friction};
\node[srv, below=\gap of ego]  (sens)
  {\textbf{Sensor actors}\\[1pt] IMU, GNSS, cameras, LiDAR};
\node[srv, below=\gap of sens] (tick)
  {\textbf{World tick}\\[1pt] synchronous, $\Delta t=0.033$\,s};
\node[srv, below=\gap of tick] (wt)
  {\textbf{Wheel telemetry}\\[1pt]
   per-wheel slip angle,\\normal load, lateral force};

\begin{scope}[on background layer]
  \node[draw, dashed, rounded corners=2pt, inner sep=3.6pt, fill=black!2,
        fit=(ego)(sens)(tick)(wt)] (carla) {};
\end{scope}
\node[colhead, above=1.4mm of carla.north]
  {CARLA server\\{\scriptsize\mdseries UE 4.26 / PhysX}};

% ======================== 2. Carla_ASU_Bridge node ========================
\node[brg, anchor=north west] (ctrl) at (5.75,0)
  {\textbf{Control application}\\[1pt]
   forwards speed and steering to\\CARLA's \emph{native} Ackermann
   controller (server-side PID)};
\node[brg, below=\gap of ctrl] (sth)
  {\textbf{Per-sensor threads}\\[1pt]
   one publish thread each;\\optional dedicated clients};
\node[brg, below=\gap of sth] (tth)
  {\textbf{Telemetry thread}\\[1pt]
   10\,Hz on its own thread --- each\\tick costs blocking CARLA RPCs};
\node[brg, below=\gap of tth] (clk)
  {\textbf{\ttfamily sensor\_clock}\\[1pt]
   every stamp drawn from the\\simulation epoch, not wall time};

\begin{scope}[on background layer]
  \node[draw, dashed, rounded corners=2pt, inner sep=3.6pt, fill=black!1,
        fit=(ctrl)(sth)(tth)(clk)] (bridge) {};
\end{scope}
\node[colhead, above=1.4mm of bridge.north]
  {\textsc{Carla\_ASU\_Bridge}\\{\scriptsize\mdseries ROS 2 lifecycle node (C\texttt{++})}};

% ============================ 3. ROS 2 topics ============================
\node[top, anchor=north west] (tsub) at (11.1,0)
  {/drive\\control/tire\_friction\\control/drag\_coefficient};
\node[note, below=1.2mm of tsub] (nsub)
  {\emph{subscribed} --- commands and\\runtime physics overrides};

\node[top, below=2.6mm of nsub] (tstate)
  {/odom\\feedback/steering\_angle\\feedback/imu\\feedback/speed};
\node[note, below=1.2mm of tstate] (nstate)
  {\emph{published} --- the signals the\\identification actually consumes};

\node[gt, below=2.6mm of nstate] (tgt)
  {feedback/tire\_forces};
\node[note, below=1.2mm of tgt] (ngt)
  {\textbf{ground truth} --- per-wheel slip\\
   angle, normal load and lateral\\
   force, straight from PhysX; no\\
   analytic tire model in between};

\begin{scope}[on background layer]
  \node[draw, dashed, rounded corners=2pt, inner sep=3.6pt, fill=black!1,
        fit=(tsub)(nsub)(tstate)(nstate)(tgt)(ngt)] (topics) {};
\end{scope}
\node[colhead, above=1.4mm of topics.north]
  {ROS 2 topic layer\\{\scriptsize\mdseries all stamped from \texttt{/clock}}};

% =========================== 4. Autonomy stack ===========================
\node[cons, anchor=north west] (pp) at (16.6,0)
  {\textbf{\ttfamily pure\_pursuit}\\[1pt]
   model-free bootstrap; needs\\only the wheelbase};
\node[cons, below=\gap of pp] (sid)
  {\textbf{\ttfamily On-Track-SysID}\\[1pt]
   residual NN $\rightarrow$ virtual sweep\\$\rightarrow$ Pacejka fit};
\node[cons, below=11mm of sid] (mpc)
  {\textbf{\ttfamily mpc\_path\_tracking}\\[1pt]
   \texttt{acados} NMPC; holds the\\admissibility gates};
\node[cons, below=\gap of mpc] (mgr)
  {\textbf{\ttfamily adaptive\_}\\\textbf{\ttfamily controller\_manager}\\[1pt]
   PP\,$\rightarrow$\,MPC arbitration};
\node[cons, below=\gap of mgr] (bch)
  {\textbf{Benchmark nodes}\\[1pt]
   1/5/10-step prediction metrics,\\$F_y$ against ground truth};

\begin{scope}[on background layer]
  \node[draw, dashed, rounded corners=2pt, inner sep=3.6pt, fill=black!2,
        fit=(pp)(sid)(mpc)(mgr)(bch)] (stack) {};
\end{scope}
\node[colhead, above=1.4mm of stack.north]
  {Autonomy stack\\{\scriptsize\mdseries this workspace}};

% ================================ arrows =================================
% CARLA <-> bridge, over the C++ client library.
% The two column groups have different heights, so carla.east and bridge.west
% do NOT share a y. Pin both endpoints to one row (the "Sensor actors" row)
% with the |- / -| intersection syntax so the arrow is exactly horizontal.
\draw[flow, <->, line width=0.7pt]
      (carla.east |- sens) -- (bridge.west |- sens)
      node[lbl, above, pos=0.5, yshift=0.3mm]{\ttfamily libcarla\_client}
      node[lbl, below, pos=0.5, yshift=-0.3mm]{RPC + streams};

% CONVENTION: every inter-column arrow begins and ends on a dashed GROUP
% boundary, never on an inner box -- an arrow that pierces the dashed frame
% reads as if it bypassed the column. The row it belongs to is carried by the
% y coordinate (|- <box>), so the arrow still lines up with its own topic.

% bridge -> published topics
\draw[flow] (bridge.east |- tstate) -- (topics.west |- tstate)
      node[lbl, above, pos=0.5]{publish};
\draw[gtflow] (bridge.east |- tgt) -- (topics.west |- tgt);

% commands travel right to left, back into the bridge
\draw[flow] (topics.west |- tsub) -- (bridge.east |- tsub)
      node[lbl, above, pos=0.5]{apply};

% topics -> autonomy stack, and the /drive return path
\draw[flow] (topics.east |- tstate) -- (stack.west |- tstate)
      node[lbl, above, pos=0.5]{state feedback};
\draw[flow] (stack.west |- tsub) -- (topics.east |- tsub)
      node[lbl, above, pos=0.5]{\ttfamily /drive};

% ground truth reaches ONLY the scoring nodes: out of the topic column at the
% tire_forces row, down the gutter, into the stack at the benchmark row.
\coordinate (gtout)  at (topics.east |- tgt);
\coordinate (gtbend) at ($(gtout)+(0.6,0)$);
\draw[gtflow] (gtout) -- (gtbend) |- (stack.west |- bch);
\node[font=\tiny, orange!80!black, rotate=90, anchor=south, inner sep=1.4pt]
      at ($(gtbend)!0.5!(gtbend |- bch)$) {validation only};

% the gated model interface, inside the stack
\draw[gateflow] (sid.south) -- (mpc.north)
      node[lbl, pos=0.5, align=center]
      {\ttfamily mpc/update\_params\\\rmfamily\emph{gated}};

% ================================= legend ================================
% Sits under columns 1--2, which are the shorter pair.
\node[anchor=north west, font=\tiny, align=left, draw, rounded corners=1.5pt,
      inner sep=4pt, fill=white, text width=88mm]
      at ($(carla.south west)+(0,-0.5)$)
  {\tikz[baseline=-0.5ex]\draw[gtflow](0,0)--(0.5,0);\ \textbf{ground-truth path.}
   Used only to \emph{score} the identifier; never inside the identification loop.\\[1.5pt]
   \tikz[baseline=-0.5ex]\draw[gateflow](0,0)--(0.5,0);\ \textbf{gated interface.}
   An identified tire set is accepted only if its grip ceiling, understeer sign and
   critical speed clear the controller's requirements (Sec.~\ref{subsec:gates}).\\[1.5pt]
   \tikz[baseline=-0.5ex]\draw[flow](0,0)--(0.5,0);\ ROS 2 topic traffic.};
\end{tikzpicture}
}
\caption{Architecture of the \bridge{} ROS~2 interface. The bridge links
against CARLA's C\texttt{++} client, drives the ego vehicle through CARLA's
native Ackermann controller, and republishes state, sensors and per-wheel
telemetry on ROS~2 topics stamped from the simulation clock. The
\texttt{feedback/\allowbreak tire\_forces} path (highlighted) carries the
simulator's own per-wheel slip angle, normal load and lateral force, and is the
ground truth against which the identifier is scored.}
\label{fig:bridge}
\end{figure*}

\begin{table}[t]
\caption{ROS~2 interface used by the identification and validation pipeline.
Namespace is \texttt{/sim} unless marked absolute.}
\label{tab:topics}
\centering
\scriptsize
\setlength{\tabcolsep}{3pt}
\begin{tabular}{@{}p{2.4cm}p{2.75cm}p{2.5cm}@{}}
\toprule
Topic & Type & Role\\
\midrule
\texttt{/odom} & \texttt{nav\_msgs/Odometry} & State $v_{x},v_{y},\omega$; twist is body-frame\\
\texttt{feedback/\allowbreak steering\_angle} & \texttt{std\_msgs/Float32} (deg) & Achieved road-wheel angle $\delta$\\
\texttt{feedback/imu} & \texttt{sensor\_msgs/Imu} & $a_{x},a_{y}$ for the friction warm-start\\
\texttt{feedback/\allowbreak tire\_forces} & \texttt{sim\_manager\_\allowbreak msgs/\allowbreak TireForces} & Per-wheel $\alpha$, $F_{z}$, $F_{y}$, $\mu$, $\omega_{w}$, $\kappa$: \textbf{ground truth}\\
\texttt{feedback/\allowbreak vehicle\_physics} & \texttt{std\_msgs/String} (JSON) & Live wheel parameters of \eqref{eq:physx}\\
\texttt{feedback/speed} & \texttt{std\_msgs/Float32} & Forward speed echo\\
\texttt{feedback/motors} & \texttt{std\_msgs/String} (JSON) & Per-wheel speed/torque/brake\\
\texttt{/drive} & \texttt{ackermann\_msgs/\allowbreak Ackermann\allowbreak Drive\allowbreak Stamped} & Control command in\\
\texttt{control/\allowbreak tire\_friction} & \texttt{std\_msgs/Float32} & Runtime grip override\\
\texttt{/clock} \emph{(abs.)} & \texttt{rosgraph\_msgs/\allowbreak Clock} & Simulation time base\\
\bottomrule
\end{tabular}
\end{table}

\paragraph{Per-wheel ground truth}
\texttt{feedback/\allowbreak tire\_forces} carries CARLA's own
\texttt{Wheel\-TelemetryData} in wheel order [FL, FR, RL, RR]. No analytic tire
model exists inside the bridge, so nothing on this topic can disagree with the
physics the vehicle is driven by. Four measured properties bound what may be
read off it. (i)~Only \texttt{lateral\_force} is force ground truth; it is the
channel every validation in this paper uses.
(ii)~\texttt{longitudinal\_force} is a friction-capacity report rather than a
measurement, tracking $\mu F_{z}$ and correlating with $ma_{x}$ at \num{0.04},
and \texttt{wheel\_torque} is an exact identity of it. (iii)~Below
\SI{0.5}{\meter\per\second} the publisher zeroes slip angle and both forces
while the normal load stays live, so the identifier gates on the all-zero
pattern as well. (iv)~The slip angle needs no sign change across the boundary,
verified rather than assumed.

Figure~\ref{fig:frames} states the conventions the rest of the paper assumes.
CARLA inherits Unreal's left-handed world while every ROS~2 topic follows
REP-103~\cite{rep103} and the identification equations are written in the body
frame with $\alpha$ and $F_{y}$ positive to the left, the SAE-style axis
convention~\cite{saej670,pacejka2012book}. The bridge is a per-channel
conversion rather than a single transform, and the trap is that the tire
telemetry's slip angle needs \emph{no} sign change, because the frame flip and
CARLA's own convention cancel: an error there is invisible in the states and
inverts the identified curve, so it was checked against a slip angle rebuilt
from the chassis state in a steady \SI{4}{\meter\per\second} turn, equal in
magnitude to within \SI{2}{\percent} front and \SI{15}{\percent} rear.

\begin{figure*}[t]
\centering
\resizebox{\textwidth}{!}{% =============================================================================
% Fig. -- Frame and sign conventions across the CARLA/ROS 2 boundary.
%
% Vector TikZ, drawn at the paper's own font sizes. Included by
% sections/platform.tex inside a figure* (full text width).
%
% Three panels, left to right:
%   (a) the two world frames and the map between them (CARLA/Unreal is
%       left-handed with Y south and yaw clockwise; ROS is ENU/REP-103).
%   (b) the vehicle body frame in which every identification quantity is
%       defined: FLU, positive steer to the left, positive lateral force to
%       the left, slip angles as in (2).
%   (c) the per-channel conversion the bridge applies, which is where a sign
%       error would enter the identification silently.
%
% Layout note: every panel's bounding box is fitted to two explicit corner
% coordinates rather than to its content, so the three panels keep the same
% height and nothing collides once the whole picture is resized to \textwidth.
% =============================================================================

\begin{tikzpicture}[
  font=\scriptsize,
  >={Stealth[length=1.7mm,width=1.5mm]},
  ax/.style      = {->, line width=0.5pt, black!75},
  car/.style     = {line width=0.7pt, black!75},
  wheel/.style   = {line width=1.5pt, black!80, line cap=round},
  frc/.style     = {->, line width=0.7pt, green!45!black},
  vel/.style     = {->, line width=0.6pt, blue!55!black},
  rot/.style     = {->, line width=0.5pt, black!60},
  lbl/.style     = {font=\tiny, inner sep=1.0pt},
  note/.style    = {font=\tiny, align=left, inner sep=1.6pt},
  hdr/.style     = {font=\scriptsize\bfseries, align=center},
  conv/.style    = {->, line width=0.6pt, black!65, dashed}
]

% ================== (a) vehicle body frame and sign rules ==================
% Single-track sketch: x forward (right), y left (up), omega CCW positive.
\begin{scope}[shift={(93mm,2mm)}]
  \coordinate (cg)  at (0,0);
  \coordinate (fr)  at (15mm,0);
  \coordinate (rr)  at (-13mm,0);

  \draw[car] (rr) -- (fr);
  \draw[wheel] ($(rr)+(-3mm,0)$) -- ($(rr)+(3mm,0)$);
  \draw[wheel, rotate around={17:(fr)}] ($(fr)+(-3mm,0)$) -- ($(fr)+(3mm,0)$);
  \fill (cg) circle (0.6mm);

  % body axes at the centre of gravity
  \draw[ax] (cg) -- ++(8mm,0) node[lbl,below] {$x$ fwd};
  \draw[ax] (cg) -- ++(0,7mm) node[lbl,left]  {$y$ left};
  \draw[rot] (-2.8mm,-2.8mm) arc (225:135:4.0mm);
  \node[lbl, anchor=north east] at (-3.2mm,-2.4mm) {$\omega>0$};

  % steer angle at the front wheel
  \draw[black!45, line width=0.35pt] (fr) -- ++(6mm,0);
  \draw[rot] ($(fr)+(5mm,0)$) arc (0:17:5mm);
  \node[lbl, anchor=west] at ($(fr)+(5.4mm,2.2mm)$) {$\delta>0$};

  % velocity and sideslip at the CG
  \draw[vel] (cg) -- ++(10mm,4.5mm) node[lbl,above] {$v$};
  \draw[rot] (5.5mm,0) arc (0:24.2:5.5mm);
  \node[lbl, anchor=west] at (6.0mm,1.2mm) {$\beta$};

  % axle lateral forces, positive to the left
  \draw[frc] (fr) -- ++(0,7mm) node[lbl,above] {$F_{y,f}>0$};
  \draw[frc] (rr) -- ++(0,7mm) node[lbl,above] {$F_{y,r}>0$};
\end{scope}

\node[note, anchor=north west, text width=62mm] at (63mm,-11mm)
  {$\alpha_{f}=\delta-\arctan\dfrac{v_{y}+l_{f}\omega}{v_{x}}$,\qquad
   $\alpha_{r}=-\arctan\dfrac{v_{y}-l_{r}\omega}{v_{x}}$\\[2pt]
   Positive $\alpha$ produces positive (leftward) $F_{y}$. Wheel order in
   every per-wheel array: [FL, FR, RL, RR].};

\coordinate (pb1) at (61mm,-27mm);
\coordinate (pb2) at (127mm,15mm);
\begin{scope}[on background layer]
  \node[draw, dashed, rounded corners=2pt, inner sep=0pt, fill=black!2,
        fit=(pb1)(pb2)] (panelb) {};
\end{scope}
\node[hdr, above=1.0mm of panelb.north] {(a) Vehicle body frame (FLU)};

% ====================== (b) per-channel conversions =========================
\node[note, anchor=north west] at (133mm,13mm)
  {\renewcommand{\arraystretch}{1.2}%
   \begin{tabular}{@{}p{17mm}p{40mm}@{}}
   \textbf{Channel} & \textbf{Bridge conversion}\\
   \hline\\[-6pt]
   \texttt{/odom} pose & $y\mapsto-y$; $\psi,\theta\mapsto-\psi,-\theta$\\
   \texttt{/odom} twist & world $\to$ body by $q^{-1}$; $v_{y}$,
   $\omega_{y}$, $\omega_{z}$ negated\\
   \texttt{feedback/imu} & $a_{y}\mapsto-a_{y}$\\
   \texttt{tire\_forces} $\alpha$ & $\deg\to\mathrm{rad}$ only: the frame flip
   and CARLA's own sign convention cancel\\
   \texttt{tire\_forces} $F_{y}$ & $F_{y}\mapsto-F_{y}$, into $+y$ left\\
   \texttt{/drive} $\delta$ & negated on the way out to CARLA\\
   \end{tabular}};

\coordinate (pc1) at (131mm,-27mm);
\coordinate (pc2) at (196mm,15mm);
\begin{scope}[on background layer]
  \node[draw, dashed, rounded corners=2pt, inner sep=0pt, fill=black!2,
        fit=(pc1)(pc2)] (panelc) {};
\end{scope}
\node[hdr, above=1.0mm of panelc.north] {(b) Conversions applied by \bridge{}};

\end{tikzpicture}
}
\caption{Frame and sign conventions across the simulator--ROS~2 boundary.
(a) All identification quantities are defined in the body frame:
$x$ forward, $y$ left, $\omega$ and $\delta$ positive counter-clockwise, and
slip angle and lateral force positive to the left, so that a positive $\alpha$
produces a positive $F_{y}$~\cite{saej670}. (b) The conversion applied per
channel. The slip angle of the tire telemetry is the one exception that needs
no sign change, because the frame flip and CARLA's own convention cancel.}
\label{fig:frames}
\end{figure*}

\section{Reference Implementation}
\label{sec:refimpl}

The identification package, the ROS~2 bridge and the controller stack used in
this paper are available as \texttt{On-Track-SysID}, extended
from~\cite{ontrack2025} (upstream:
\url{https://github.com/ForzaETH/On-Track-SysID}), \bridge{}
(\url{https://github.com/ebrahimabdelghfar/Carla_ASU_Bridge})%
~\cite{carlaasubridge}, and the accompanying MPC and benchmarking packages,
released together with the scenario definitions, the exported result CSVs and
the figure and sweep generators at
\url{https://github.com/ebrahimabdelghfar/Ebrahim_Master_Thesis_repo/tree/feat-follow-premade-racing-track}.
